\section{Transducción sísmica: alternativas y restricción impuesta por el SM-24}
\label{sec:transductores}

El transductor fija la componente observada, sensibilidad, banda y ruido antes de que
intervenga el amplificador. La electrónica puede amplificar, filtrar o compensar
parcialmente su respuesta, pero no cambia la dinámica mecánica. Esta sección compara las
familias aplicables a Rayleigh y cuantifica el desajuste entre el SM-24 disponible, de
\SI{10}{\hertz}, y la banda requerida para investigar hasta \SI{50}{\meter}.

\subsection{Qué debe preservar el transductor}

Una onda Rayleigh produce movimiento elíptico en el plano vertical; un MASW activo
convencional usa receptores verticales para captar su componente dominante. Los sensores
horizontales o triaxiales son necesarios para Love, elipticidad, HVSR o ciertos arreglos
pasivos, no para una curva Rayleigh básica~\cite{Foti2014,Park1999}.

El registro de un canal puede representarse como

\begin{equation}
\begin{aligned}
    y_j(t)&=h_{s,j}(t)*q_j(t)+n_j(t),\\
    Y_j(f)&=H_{s,j}(f)Q_j(f)+N_j(f),
\end{aligned}
    \label{eq:transductor-medicion}
\end{equation}

donde $q_j$ es desplazamiento, velocidad o aceleración; $H_{s,j}$ incluye sensor y
acoplamiento; y $N_j$ es ruido. MASW usa diferencias de fase: una respuesta idéntica en
todos los canales se cancela, mientras diferencias de amortiguamiento, inclinación o
acoplamiento introducen fase diferencial y sesgan la velocidad aparente.

Para obtener la dispersión importa principalmente que cada frecuencia conserve una
relación de fase coherente entre canales y una relación señal--ruido útil. La variable
física medida es secundaria mientras la respuesta del sensor sea conocida y uniforme
en todo el arreglo.

\subsection{Alternativas y costos de referencia}

La comparación reúne magnitud, banda, ruido, alimentación, ejes y precio del elemento;
no incluye cableado, carcasa, PCB, digitalización, importación ni impuestos. Los precios
se consultaron el 18 de agosto de 2026 y sólo indican órdenes de magnitud.

\begin{table*}[ht]
\centering
\footnotesize
\caption{Alternativas de transducción y costo del elemento sensor.}
\begin{tabularx}{\textwidth}{@{}>{\raggedright\arraybackslash}p{0.15\textwidth} >{\raggedright\arraybackslash}p{0.23\textwidth} >{\raggedright\arraybackslash}p{0.19\textwidth} Y@{}}
\toprule
\textbf{Referencia} & \textbf{Magnitud y banda} & \textbf{Precio publicado} &
\textbf{Evaluación para el arreglo} \\
\midrule
SM-24, \SI{10}{\hertz} & Geófono vertical; velocidad; respuesta plana sobre $f_n$ &
69,95 USD con disco aislante & Pasivo y estándar en MASW; 24 elementos cuestan
aproximadamente 1.679 USD, pero la respuesta cae bajo $f_n$~\cite{SparkFunSM24} \\
RTC, \SI{4.5}{\hertz} & Geófono vertical; velocidad; mayor extensión hacia bajas
frecuencias & 92--154 USD según carcasa & Sustitución directa: 2.208--3.696 USD para 24
puntos~\cite{RTClark45} \\
ADXL355 & Aceleración desde continua; triaxial, ADC de 20 bits & 70,93 USD por unidad;
41,84 USD a 1.000 & Integración simple, pero deben validarse ruido, sincronización y
acoplamiento del nodo completo~\cite{DigiKeyADXL355,AnalogDevicesADXL355} \\
PCB 393B31 & Aceleración, \SIrange{0.1}{200}{\hertz}; $1~\mu g$ RMS & Por cotización &
Muy bajo ruido, pero requiere IEPE de \SIrange{24}{28}{\volt} y adquisición compatible
~\cite{PCB393B31} \\
Güralp 3ESPC & Velocidad triaxial, 60 s--\SI{100}{\hertz} & Por cotización & Cubre bajas
frecuencias con bajo ruido; es un instrumento profesional completo~\cite{Guralp3ESPC} \\
\bottomrule
\end{tabularx}
\end{table*}

Los geófonos siguen siendo la opción natural para MASW activo por sensibilidad,
robustez y pasividad. Reducir $f_n$ conserva mejor las longitudes de onda grandes, pero
aumenta el costo. Un MEMS tampoco es automáticamente más económico: en el ADXL355 la
densidad de ruido especificada es $22{,}5~\mu g/\sqrt{\mathrm{Hz}}$, y el costo real
depende además de encapsulado, calibración y comunicación~\cite{AnalogDevicesADXL355}.

\subsection{SM-24: restricción de partida}

El SM-24 de \SI{10}{\hertz} se empleó porque estaba disponible en la Facultad, no porque
resultara de optimizar la profundidad. Su $f_n$ dificulta recuperar los pocos hertz
asociados a \SI{50}{\meter}; aun así, es vertical, pasivo, robusto, documentado y de
costo compatible con un arreglo académico. Su salida analógica permite además estudiar
el acondicionamiento como parte central del proyecto.

\begin{figure}[!ht]
\centering
\includegraphics[width=\columnwidth]{sm24_esquema_interno.png}\par\smallskip
\includegraphics[width=0.80\columnwidth]{sm24_fotografia_cableado.png}
\caption{Principio constructivo de un geófono electromagnético y SM-24 encapsulado
con cable y punta de acoplamiento al terreno.}
\label{fig:sm24-fisico-respuesta}
\end{figure}

Sus parámetros nominales principales son~\cite{SM24}:

\begin{table}[!ht]
\centering
\scriptsize
\caption{Parámetros nominales del SM-24.}
\begin{tabularx}{\columnwidth}{@{}Y>{\raggedleft\arraybackslash}p{0.36\columnwidth}@{}}
\toprule
\textbf{Parámetro} & \textbf{Valor} \\
\midrule
Frecuencia natural & \SI{10}{\hertz} $\pm2{,}5\%$ \\
Sensibilidad & \SI{28.8}{\volt\second\per\meter} $\pm2{,}5\%$ \\
Amortiguamiento abierto & $\zeta=0{,}25$ \\
Amortiguamiento con \SI{1339}{\ohm} & $\zeta=0{,}60$ \\
Resistencia de bobina & \SI{375}{\ohm} $\pm2{,}5\%$ \\
Masa móvil & \SI{11}{\gram} \\
Frecuencia espuria típica & $>\SI{240}{\hertz}$ \\
\bottomrule
\end{tabularx}
\end{table}

\subsection{Modelo dinámico del SM-24: velocidad y aceleración}

Sea $x_g$ el desplazamiento de la carcasa unida al suelo, $x_m$ el de la masa móvil y
$z=x_m-x_g$ el desplazamiento relativo. El sistema masa--resorte--amortiguador satisface

\begin{equation}
\begin{aligned}
    m\ddot z+c\dot z+kz&=-m\ddot x_g,\\
    \omega_n&=\sqrt{\frac{k}{m}},
    &\zeta&=\frac{c}{2\sqrt{km}}.
\end{aligned}
    \label{eq:sm24-mecanico}
\end{equation}

La bobina genera una fuerza electromotriz proporcional a la velocidad relativa:

\begin{equation}
    v_o(t)=G_0\dot z(t),
\end{equation}

por lo que, tomando como entrada la velocidad del suelo,

\begin{equation}
    H_v(s)=\frac{V_o(s)}{\dot X_g(s)}
    =-G_0\frac{s^2}{s^2+2\zeta\omega_ns+\omega_n^2}.
    \label{eq:sm24-velocidad}
\end{equation}

Con $r=\omega/\omega_n$, la magnitud normalizada es

\begin{equation}
    \frac{|H_v(j\omega)|}{G_0}
    =\frac{r^2}{\sqrt{(1-r^2)^2+(2\zeta r)^2}}.
    \label{eq:sm24-magnitud}
\end{equation}

La aceleración del suelo es $A_g(s)=s^2X_g(s)=s\dot X_g(s)$. Por tanto, la misma salida
del SM-24 puede expresarse respecto de la aceleración como

\begin{equation}
    H_a(s)=\frac{V_o(s)}{A_g(s)}
    =\frac{H_v(s)}{s}
    =-G_0\frac{s}{s^2+2\zeta\omega_ns+\omega_n^2}.
    \label{eq:sm24-aceleracion}
\end{equation}

Su magnitud, normalizada por $G_0/\omega_n$, es

\begin{equation}
    \frac{\omega_n|H_a(j\omega)|}{G_0}
    =\frac{r}{\sqrt{(1-r^2)^2+(2\zeta r)^2}}.
    \label{eq:sm24-magnitud-aceleracion}
\end{equation}

El geófono sin compensar se usa habitualmente referido a velocidad porque
$H_v\rightarrow-G_0$ por encima de $f_n$: allí su respuesta es aproximadamente plana.
Referido a aceleración, $H_a$ es pasabanda y no presenta una meseta con el
amortiguamiento nominal.

Por debajo de $f_n$, $|H_v|$ cae como $r^2$ (\SI{40}{\decibel\per decade}) y $|H_a|$
como $r$ (\SI{20}{\decibel\per decade}); son dos referencias de una misma tensión de
salida, no dos modos de funcionamiento distintos.

\begin{figure}[!ht]
\centering
\includegraphics[width=\columnwidth]{respuesta_sm24_vel_acel.png}
\caption{Modelo nominal del SM-24 en circuito abierto. Ambas referencias se muestran
en los mismos ejes: bajo $f_n$ caen \SI{40}{\decibel\per decade} y
\SI{20}{\decibel\per decade}; sobre $f_n$, la referencia a velocidad se aplana y la
referencia a aceleración cae \SI{20}{\decibel\per decade}.}
\label{fig:sm24-modelo}
\end{figure}

\subsection{La banda exigida por el objetivo de profundidad}

La propuesta original planteó caracterizar hasta \SI{50}{\meter}. Como estimación de
diseño se usa

\begin{equation}
    \lambda_R=\frac{c_R}{f}.
\end{equation}

Longitudes de \SIrange{100}{150}{\meter} y velocidades Rayleigh de
\SIrange{150}{500}{\meter\per\second} sitúan la información profunda en
\SIrange{1}{5}{\hertz}. No es una conversión universal profundidad--frecuencia, sino una
comparación preliminar con la respuesta del sensor.

El conflicto queda cuantificado al evaluar la ecuación~\eqref{eq:sm24-magnitud}:

\begin{table}[!ht]
\centering
\scriptsize
\caption{Respuesta nominal del SM-24 en circuito abierto, normalizada respecto de
velocidad y aceleración.}
\begin{tabular}{@{}rrrr@{}}
\toprule
$f$ & $f/f_n$ & $|H_v|/G_0$ & $\omega_n|H_a|/G_0$ \\
\midrule
\SI{1}{\hertz}  & 0,10 & $-40{,}0$ dB & $-19{,}9$ dB \\
\SI{2}{\hertz}  & 0,20 & $-27{,}6$ dB & $-13{,}7$ dB \\
\SI{5}{\hertz}  & 0,50 & $-10{,}0$ dB & $-4{,}0$ dB \\
\SI{10}{\hertz} & 1,00 & $+6{,}0$ dB & $+6{,}0$ dB \\
\SI{20}{\hertz} & 2,00 & $+2{,}0$ dB & $-4{,}0$ dB \\
\bottomrule
\end{tabular}
\end{table}

En la referencia a velocidad, la salida cae \SI{27.6}{\decibel} a \SI{2}{\hertz} y
\SI{40}{\decibel} a \SI{1}{\hertz}; normalizada respecto de aceleración, cae
\SI{13.7}{\decibel} y \SI{19.9}{\decibel}. La diferencia refleja pendientes de
\SI{40}{\decibel\per decade} y \SI{20}{\decibel\per decade}; no crea señal adicional
en los terminales.

\subsection{Implicación para el acondicionamiento}

La tesis asume el desajuste del SM-24 como problema de instrumentación: determinar cuánto
contenido de baja frecuencia puede conservarse mediante una cadena analógica de bajo
ruido que compense la región atenuada y preserve fase y rango dinámico. No se considera
al SM-24 óptimo ni se atribuye a la electrónica la capacidad de eliminar su límite
mecánico. La siguiente sección desarrolla esa estrategia.
